\section{Related Work}
\label{sec:related}

\paragraph{Attack Investigation and Triage} Provenance-based systems build causal graphs of host activity and make them tractable with rules or learned models. NoDoze~\citep{nodoze} ranks alerts by the anomaly of their surrounding provenance to combat alert fatigue, and Kairos~\citep{kairos} detects and reconstructs intrusions from whole-system provenance. These systems decide what an analyst should look at, not which query to run next. Clouseau~\citep{clouseau} reconstructs attack narratives from incident logs with a hierarchy of LLM agents that choose their own queries, and reports that open-weight models trail a proprietary one on complex scenarios; its authors note that narrative reconstruction exceeds what triage, which seeks a quick accept-or-escalate decision, requires. \hesp{} targets that decision, with small local models, and moves the procedure out of the model. ExCyTIn-Bench~\citep{excytin} and GUIDE~\citep{guide} provide data for LLM and ML investigation; \S\ref{sec:realdata} measures why neither can test an investigation policy.

\paragraph{LLMs for Security} LLM agents have been applied to penetration testing~\citep{pentestgpt} and capture-the-flag challenges~\citep{cybench}. Work on securing agents constrains what an agent may do: IsolateGPT~\citep{isolategpt} isolates the execution of LLM apps, and Progent~\citep{progent} enforces privilege policies on tool calls. \hesp{} shares their premise that an agent's actions should be governed outside the model, but it governs the investigation procedure (which probe runs, which verdict is accepted, when to stop) rather than access rights, and it measures the effect on task success.

\paragraph{Tool-Using Agents} ReAct~\citep{react} interleaves reasoning and actions, Reflexion~\citep{reflexion} adds verbal self-feedback, and Toolformer~\citep{toolformer} learns when to call tools. In these designs the model decides what to do next and when to finish. \hesp{} moves probe selection, and optionally stopping, into an explicit controller, and uses the model's own choices as a baseline configuration.

\paragraph{Sequential Diagnosis} Choosing tests by expected information gain goes back to Lindley~\citep{lindley1956} and Bayesian experimental design~\citep{chaloner1995}. It underlies classic model-based diagnosis~\citep{dekleer1987} and has near-optimality guarantees for greedy active learning and noisy hypothesis identification~\citep{dasgupta2004,golovin2010}. We use the textbook criterion unchanged; our question is empirical, namely how such a criterion interacts with LLM planners of different strengths. Language models can be reasonably calibrated on some self-assessment tasks~\citep{kadavath2022}, but our elicited tables show that this does not extend to the probe-outcome likelihoods a diagnosis controller needs.

\paragraph{Evaluation Methodology} Pre-registration separates confirmatory from exploratory analysis~\citep{nosek2018}, and pitfalls in ML-based security evaluation are well documented~\citep{arp2022dos}. We pre-register every study, freeze source and table hashes, and publish errata.
